\documentclass[11pt,a4paper]{article}
\usepackage{fontspec,amsmath,amssymb,amsthm,microtype}
\setmainfont{lmroman10-regular.otf}[BoldFont=lmroman10-bold.otf,ItalicFont=lmroman10-italic.otf,BoldItalicFont=lmroman10-bolditalic.otf]
\usepackage{xcolor}
\usepackage[most]{tcolorbox}
\usepackage{mathrsfs}
\providecommand{\lt}{<}
\providecommand{\gt}{>}
\setmonofont{DejaVuSansMono.ttf}[BoldFont=DejaVuSansMono-Bold.ttf,ItalicFont=DejaVuSansMono-Oblique.ttf,BoldItalicFont=DejaVuSansMono-BoldOblique.ttf,Scale=0.86]
\newfontfamily\leanfallback{STIXTwoMath-Regular.otf}[Scale=0.86]
\definecolor{leanrule}{gray}{0.62}
\definecolor{leanback}{gray}{0.962}
\definecolor{leantitle}{gray}{0.22}
\definecolor{leanmuted}{gray}{0.38}
\newtcolorbox{leanbox}[1]{enhanced,breakable,colback=leanback,colframe=leanrule,
  boxrule=0.4pt,arc=0.6mm,left=2.4mm,right=2.4mm,top=1.8mm,bottom=1.8mm,
  before skip=0.8em,after skip=1em,fonttitle=\small\bfseries,coltitle=leantitle,
  colbacktitle=leanback,titlerule=0.3pt,title={#1}}
\newlength{\leancharwidth}
\newenvironment{leancode}{\par\smallskip\begingroup\ttfamily\footnotesize\raggedright
  \settowidth{\leancharwidth}{M}\setlength{\parskip}{0pt}\setlength{\parindent}{0pt}
  \frenchspacing}{\par\endgroup\smallskip}
\newcommand{\leanline}[2]{\hangindent=\dimexpr#1\leancharwidth+4\leancharwidth\relax
  \hangafter=1\noindent\hspace*{#1\leancharwidth}#2\par}
\newcommand{\leansource}[1]{\par{\footnotesize\color{leanmuted}#1}\par}
\newcommand{\leanlabel}[1]{\par\smallskip\noindent{\small\bfseries #1}\par}
\usepackage[margin=25mm]{geometry}
\usepackage[unicode,hidelinks]{hyperref}
\hypersetup{pdfauthor={LangevinLean Project}}
\setlength{\emergencystretch}{3em}
\setlength{\parskip}{0.4em}
\setlength{\parindent}{0pt}
\allowdisplaybreaks[2]
\raggedbottom
\title{\texorpdfstring{An entropy upper bound at the initial time}{An entropy upper bound at the initial time}}
\author{LangevinLean Project}
\date{}
\begin{document}\maketitle
Draft mathematical exposition; these notes may contain mistakes.


Let \(H_0<\infty\) be the initial relative entropy and \(G_0=\int\|\nabla U\|^2d\mu\), finite by the preceding entropy-moment estimate. With \(\delta(t)=1-4tL\), define

\[B(t)=H_0-\frac d2\log\delta(t)+\frac{t}{4\delta(t)}\left[2(1+L^2t^2)G_0+4tdL^2\right].\]

The updated Gibbs density comparison, data processing, and propagated gradient bound give finite interpolation entropy \(H_t\le B(t)\) for sufficiently small positive \(t\). The denominator tends to one, so \(B\) is continuous at zero and \(B(0)=H_0\). This upper bound supplies the initial value in the scalar comparison. No initial Fisher information, position second moment, or continuity assumption on the entropy path at zero is needed.


\section{\texorpdfstring{The explicit initial-endpoint majorant}{The explicit initial-endpoint majorant}}


Put \(\delta(t)=1-4tL\) and \(G_0=\mathbb E_\mu\|\nabla U\|^2\). Define the explicit upper bound \(B(t)=H_0-\frac d2\log\delta(t)+\frac{t}{4\delta(t)}[2(1+L^2t^2)G_0+4tdL^2]\). The logarithmic term is the normalizer correction from the updated Gibbs density; the final term uses the propagated gradient moment. The definition is a scalar function at all times, while its entropy interpretation below requires \(t>0\) and \(\delta(t)>0\).



The majorant combines the data-processing entropy with the Gaussian normalizer correction and a gradient-energy correction. Its explicit denominator is positive for all sufficiently small positive times, so the expression is an ordinary finite scalar there.

\begin{leanbox}{Lean: \texttt{interpolationEntropyMajorant}}
\leansource{\texttt{interpolationEntropyMajorant} \quad UnadjustedLangevin/\allowbreak{}LSI/\allowbreak{}InitialEntropy.lean, lines 25–32}
\begin{leancode}
\leanline{0}{def\ interpolationEntropyMajorant\ (U\ :\ E\ →\ ℝ)\ (L\ :\ ℝ≥0)\ (μ\ :\ Measure\ E)\ (t\ :\ ℝ)\ :\ ℝ\ :=}
\leanline{2}{(InformationTheory.klDiv\ μ\ (ULA.gibbsMeasure\ U)).toReal\ +}
\leanline{4}{-((Module.finrank\ ℝ\ E\ :\ ℝ)\ /\ 2)\ *\ Real.log\ (1\ -\ 4\ *\ t\ *\ (L\ :\ ℝ))\ +}
\leanline{4}{(t\ /\ (1\ -\ 4\ *\ t\ *\ (L\ :\ ℝ))\ /\ 4)\ *}
\leanline{6}{(2\ *\ (1\ +\ (L\ :\ ℝ)\ \textasciicircum{}\ 2\ *\ t\ \textasciicircum{}\ 2)\ *\ (∫\ x,\ ‖gradient\ U\ x‖\ \textasciicircum{}\ 2\ ∂μ)\ +}
\leanline{8}{4\ *\ t\ *\ (Module.finrank\ ℝ\ E\ :\ ℝ)\ *\ (L\ :\ ℝ)\ \textasciicircum{}\ 2)}
\leanline{0}{\strut}
\leanline{0}{@[simp]}
\end{leancode}
\leanlabel{How the Lean proof follows the argument}
\texttt{interpolation\allowbreak{}Entropy\allowbreak{}Majorant} is a definition with defining expression \texttt{(Information\allowbreak{}Theory.\allowbreak{}kl\allowbreak{}Div \ensuremath{\mu} (ULA.\allowbreak{}gibbs\allowbreak{}Measure U)).\allowbreak{}to\allowbreak{}Real +}. The right-hand side specifies the object used in the mathematical statement.
\end{leanbox}


Earlier results used here: \cite[Section~3.26.3]{LangevinLean2026ULASC}.

\section{\texorpdfstring{The majorant at zero}{The majorant at zero}}


At \(t=0\), the logarithm has argument \(\delta(0)=1\), hence is zero. The coefficient \(t/(4\delta(t))\) is zero as well. Both correction terms vanish exactly, proving \(B(0)=H_0\). This is a direct algebraic evaluation of the constructed scalar function. No limit of the entropy path is assumed.



At zero the logarithmic term is \(\log1=0\), and the entire energy correction has a factor \(t\). Hence the value is precisely the initial entropy, even when the initial law has no smooth density.

\begin{leanbox}{Lean: \texttt{interpolationEntropyMajorant\_\allowbreak{}zero}}
\leansource{\texttt{interpolationEntropyMajorant\_\allowbreak{}zero} \quad UnadjustedLangevin/\allowbreak{}LSI/\allowbreak{}InitialEntropy.lean, lines 33–36}
\begin{leancode}
\leanline{0}{theorem\ interpolationEntropyMajorant\_zero\ (U\ :\ E\ →\ ℝ)\ (L\ :\ ℝ≥0)\ (μ\ :\ Measure\ E)\ :}
\leanline{4}{interpolationEntropyMajorant\ U\ L\ μ\ 0\ =}
\leanline{6}{(InformationTheory.klDiv\ μ\ (ULA.gibbsMeasure\ U)).toReal}
\end{leancode}
\leanlabel{How the Lean proof follows the argument}
\texttt{interpolation\allowbreak{}Entropy\allowbreak{}Majorant\_\allowbreak{}zero}: the concluding step is \texttt{simp [interpolation\allowbreak{}Entropy\allowbreak{}Majorant]}.
\end{leanbox}


Earlier results used here: \cite[Section~3.26.3]{LangevinLean2026ULASC}.

\section{\texorpdfstring{Continuity of the majorant at zero}{Continuity of the majorant at zero}}


The denominator \(\delta(t)\) has value one at zero, so it stays nonzero locally. The logarithm is continuous at that same positive value. All other terms in \(B\) are polynomial or constant scalar operations. Their continuity proves \(B(t)\to B(0)=H_0\) as \(t\to0\). The initial gradient integral is an already defined scalar; its finiteness in the entropy application is established by the finite-KL moment theorem.



The denominator tends to one and is nonzero near zero. Continuity of the logarithm at one, polynomial factors, and reciprocal then proves continuity of the full majorant at the endpoint.

\begin{leanbox}{Lean: \texttt{interpolationEntropyMajorant\_\allowbreak{}continuousAt\_\allowbreak{}zero}}
\leansource{\texttt{interpolationEntropyMajorant\_\allowbreak{}continuousAt\_\allowbreak{}zero} \quad UnadjustedLangevin/\allowbreak{}LSI/\allowbreak{}InitialEntropy.lean, lines 40–44}
\begin{leancode}
\leanline{0}{theorem\ interpolationEntropyMajorant\_continuousAt\_zero}
\leanline{4}{(U\ :\ E\ →\ ℝ)\ (L\ :\ ℝ≥0)\ (μ\ :\ Measure\ E)\ :}
\leanline{4}{ContinuousAt\ (interpolationEntropyMajorant\ U\ L\ μ)\ 0}
\end{leancode}
\leanlabel{How the Lean proof follows the argument}
\texttt{interpolation\allowbreak{}Entropy\allowbreak{}Majorant\_\allowbreak{}continuous\allowbreak{}At\_\allowbreak{}zero}: the proof is the application \texttt{unfold interpolation\allowbreak{}Entropy\allowbreak{}Majorant}. It specializes that result with the hypotheses and objects in the exact statement.
\end{leanbox}

\section{\texorpdfstring{The interpolation entropy lies below the upper bound}{The interpolation entropy lies below the upper bound}}


For positive \(t\) with \(\delta=1-4tL>0\), choose the comparison Gaussian time \(s=t/\delta\). The updated Gibbs law has an explicit relative density bounded by \((s/t)^{d/2}\exp(s\|\nabla U(y)\|^2/4)\). Since \(s/t=\delta^{-1}\), this is \(\exp[-(d/2)\log\delta+(s/4)\|\nabla U(y)\|^2]\).



The relative density identifies the Radon-Nikodym derivative and proves absolute continuity of the updated Gibbs reference law. The change-of-reference entropy estimate then gives finite \(H_t\) and \(H_t\le H_0-(d/2)\log\delta+(s/4)G_t\). The coefficient \(s/4\) is nonnegative, so inserting \(G_t\le2(1+L^2t^2)G_0+4tdL^2\) proves \(H_t\le B(t)\).



The density-ratio comparison bounds output entropy by the majorant whenever \(1-4tL>0\). The finite score moment supplies all integrability required for the likelihood decomposition.

\begin{leanbox}{Lean: \texttt{interpolationLaw\_\allowbreak{}entropy\_\allowbreak{}le\_\allowbreak{}majorant}}
\leansource{\texttt{interpolationLaw\_\allowbreak{}entropy\_\allowbreak{}le\_\allowbreak{}majorant} \quad UnadjustedLangevin/\allowbreak{}LSI/\allowbreak{}InitialEntropy.lean, lines 48–100}
\begin{leancode}
\leanline{0}{theorem\ interpolationLaw\_entropy\_le\_majorant}
\leanline{4}{(U\ :\ E\ →\ ℝ)\ (L\ :\ ℝ≥0)\ (hL\ :\ 0\ <\ L)\ (hU\ :\ Differentiable\ ℝ\ U)}
\leanline{4}{(hLip\ :\ LipschitzWith\ L\ (gradient\ U))}
\leanline{4}{[IsProbabilityMeasure\ (ULA.gibbsMeasure\ U)]}
\leanline{4}{(μ\ :\ Measure\ E)\ [IsProbabilityMeasure\ μ]}
\leanline{4}{(hKL\ :\ InformationTheory.klDiv\ μ\ (ULA.gibbsMeasure\ U)\ ≠\ ⊤)}
\leanline{4}{(t\ :\ ℝ)\ (ht\ :\ 0\ <\ t)\ (hδ\ :\ 0\ <\ 1\ -\ 4\ *\ t\ *\ (L\ :\ ℝ))\ :}
\leanline{4}{InformationTheory.klDiv\ (interpolationLaw\ U\ μ\ t)\ (ULA.gibbsMeasure\ U)\ ≠\ ⊤\ ∧}
\leanline{6}{(InformationTheory.klDiv\ (interpolationLaw\ U\ μ\ t)\ (ULA.gibbsMeasure\ U)).toReal\ ≤}
\leanline{8}{interpolationEntropyMajorant\ U\ L\ μ\ t}
\end{leancode}
\leanlabel{How the Lean proof follows the argument}
\texttt{interpolation\allowbreak{}Law\_\allowbreak{}entropy\_\allowbreak{}le\_\allowbreak{}majorant}: the intermediate facts include \texttt{\ensuremath{\delta}}, \texttt{s}, \texttt{b}, \texttt{hratio}, \texttt{hmult}, in source order; the source steps apply or rewrite \texttt{div\_\allowbreak{}mul\_\allowbreak{}cancel\ensuremath{{}_{0}}}, \texttt{h\ensuremath{\delta}.\allowbreak{}ne'}, \texttt{Real.\allowbreak{}exp\_\allowbreak{}add}, \texttt{Real.\allowbreak{}log\_\allowbreak{}inv}, \texttt{hb.\allowbreak{}trans}. The remaining scalar rearrangements are checked by \texttt{field\_\allowbreak{}simp}, \texttt{ring}, \texttt{linarith}.
\end{leanbox}


Earlier results used here: \cite[Section~3.26.3]{LangevinLean2026ULASC}.

\section{\texorpdfstring{The entropy majorant on a positive neighborhood}{The entropy majorant on a positive neighborhood}}


For all sufficiently small positive \(t\), specifically \(t<1/(8L)\), one has \(1-4tL>0\). The previous positive-time entropy bound therefore applies eventually in the right neighborhood of zero, giving \(H_t\le B(t)\). Combined with continuity of \(B\) and \(B(0)=H_0\), this supplies the upper initial boundary used in scalar comparison. It does not insert an assumption that \(H_t\) is differentiable or continuous at zero.



A sufficiently small positive neighborhood satisfies the denominator condition. Thus the entropy upper bound is available along every approach to zero through positive times, exactly the endpoint control used in scalar comparison.

\begin{leanbox}{Lean: \texttt{interpolationLaw\_\allowbreak{}entropy\_\allowbreak{}eventually\_\allowbreak{}le\_\allowbreak{}majorant}}
\leansource{\texttt{interpolationLaw\_\allowbreak{}entropy\_\allowbreak{}eventually\_\allowbreak{}le\_\allowbreak{}majorant} \quad UnadjustedLangevin/\allowbreak{}LSI/\allowbreak{}InitialEntropy.lean, lines 104–121}
\begin{leancode}
\leanline{0}{theorem\ interpolationLaw\_entropy\_eventually\_le\_majorant}
\leanline{4}{(U\ :\ E\ →\ ℝ)\ (L\ :\ ℝ≥0)\ (hL\ :\ 0\ <\ L)\ (hU\ :\ Differentiable\ ℝ\ U)}
\leanline{4}{(hLip\ :\ LipschitzWith\ L\ (gradient\ U))}
\leanline{4}{[IsProbabilityMeasure\ (ULA.gibbsMeasure\ U)]}
\leanline{4}{(μ\ :\ Measure\ E)\ [IsProbabilityMeasure\ μ]}
\leanline{4}{(hKL\ :\ InformationTheory.klDiv\ μ\ (ULA.gibbsMeasure\ U)\ ≠\ ⊤)\ :}
\leanline{4}{∀ᶠ\ t\ in\ {\leanfallback 𝓝}[>]\ (0\ :\ ℝ),}
\leanline{6}{(InformationTheory.klDiv\ (interpolationLaw\ U\ μ\ t)\ (ULA.gibbsMeasure\ U)).toReal\ ≤}
\leanline{8}{interpolationEntropyMajorant\ U\ L\ μ\ t}
\end{leancode}
\leanlabel{How the Lean proof follows the argument}
\texttt{gibbs\_\allowbreak{}frozen\_\allowbreak{}exponent\_\allowbreak{}le} starts with \texttt{hmodel}, the smooth upper model, and \texttt{hdiff}, the gradient-difference bound. \texttt{hinner} controls the resulting mixed product. After \texttt{hx} and \texttt{hy} rewrite both squares in a common displacement, the positive factor involving the two times clears denominators and \texttt{hscale'} cancels the remaining quadratic coefficient. \texttt{interpolation\allowbreak{}Law\_\allowbreak{}entropy\_\allowbreak{}eventually\_\allowbreak{}le\_\allowbreak{}majorant} restricts to \texttt{hsmall}, sufficiently small positive times, and obtains \texttt{h\ensuremath{\delta}} before applying the finite majorant theorem. \texttt{affine\_\allowbreak{}decay\_\allowbreak{}of\_\allowbreak{}local\_\allowbreak{}ae\_\allowbreak{}deriv\_\allowbreak{}le} compares on a positive starting interval and defines \texttt{F} with the continuous majorant in place of the initial entropy value. Its \texttt{hlim} and \texttt{le\_\allowbreak{}of\_\allowbreak{}tendsto\_\allowbreak{}of\_\allowbreak{}tendsto} pass that comparison to the zero-time limit.
\end{leanbox}


Earlier results used here: \cite[Section~3.26.3]{LangevinLean2026ULASC}.

\begin{thebibliography}{ULA-LSI}
\bibitem[ULA-SC]{LangevinLean2026ULASC}
LangevinLean Project. \emph{Unadjusted Langevin under Strong Convexity}. LL-01, 2026. Draft report, content revision d50c0f64815b. \url{https://odakyildiz.github.io/LangevinLean/docs/pdf/UnadjustedLangevin/FullProof.pdf}.
\end{thebibliography}

\end{document}
